\section{产品经理路径：从结构化数据到非结构化生成}

本节先把人物经历转成方法论线索。王冠的职业路径不是背景介绍而已，它解释了他为什么会从“如何提供算法能力”一路走到“如何设计模型能力和生成系统”。

本章先看王冠的职业路径。它横跨大数据、CV/深度学习、预训练模型、Moonshot 模型产品经理和 ONE2X 创业。这个路径的教学意义在于，它提供了一个产品经理视角的 AI 技术史：AI 不只是算法模型的演进，也是数据表达方式、产品位置和用户感知方式的演进。

\begin{figure}[H]
\centering
\includegraphics[width=0.96\textwidth]{product-manager-trajectory.png}
\caption{王冠的 AI 产品路径：大数据、CV/深度学习、预训练、Moonshot 和 ONE2X。自制概念图，依据 00:02:00--00:28:39 对谈内容整理。}
\end{figure}

\begin{knowledgebox}{读图：路径的共同线索是“降低算法能力使用门槛”}
从百度用户画像，到算法开放平台、开源框架、算法生产力工具，再到 Moonshot 模型产品，王冠的工作一直围绕“让用户更低门槛获得某种算法能力”。ONE2X 只是这条线在视频生成和生成系统上的延伸。
\end{knowledgebox}

\subsection{AI 的三个产品时代}

本节从经历里抽出第一个判断：AI 周期的变化，关键不只是模型大小，而是模型能够拟合什么形态的数据。

王冠没有用 1.0/2.0/3.0，而是用数据是否结构化来划分 AI 周期。传统机器学习和早期 CV 依赖结构化标签：用户画像、矩形框、类别标签、明确字段。大模型则拟合非结构化数据：语言、图像、视频和连续世界表达。非结构化数据的表达力更强，也更接近真实世界的连续性。

\begin{knowledgebox}{术语消化：结构化与非结构化数据}
\begin{tabularx}{\textwidth}{p{0.22\textwidth}p{0.34\textwidth}X}
\toprule
术语 & 一句话解释 & 本期关系 \\
\midrule
结构化数据 & 有明确字段、标签、坐标或类别的数据 & 传统机器学习和 CV 时代的核心输入。 \\
非结构化数据 & 文本、图像、视频等连续表达，不预先压成固定字段 & 大模型和生成系统的核心训练对象。 \\
User profile & 用户画像，用标签描述用户偏好和行为 & 推荐系统中用户侧建模的基础。 \\
模型产品经理 & 设计模型效果、能力和评测，使用户能感知到能力 & Moonshot 经历中形成的关键岗位。 \\
\bottomrule
\end{tabularx}
\end{knowledgebox}

\subsection{产品经理为什么没有被模型替代}

王冠认为 AI 时代产品经理的价值变大了，而不是变小了。原因有两层。第一，模型像人的 System1：它把大量信息提前压缩进系统里，能快速、本能地对输入反应。产品经理要设计这个 System1 应该具备什么能力，并让这些能力被训练出来、被用户感知。第二，模型输出取决于前面有哪些 token；产品经理要设计 workflow、agent 框架、领域知识库和 context，让模型在生成时获得更多有效 token。

\begin{importantbox}{AI 产品经理的两项新工作}
第一，定义模型能力：什么效果值得训练，怎样评测，用户如何感知。第二，设计 context：怎样把业务 know-how、工作流、工具和领域知识转成模型可利用的有效 token。
\end{importantbox}

\subsection{本章小结}

王冠的路径说明，AI 产品经理不是被模型能力吞没，而是被推到更前台。模型能力越强，越需要有人定义能力、组织 context、设计任务形态，并把能力变成用户能理解、能付费、能持续使用的产品。

